\begin{scanpage}{0173}{155}
there are a number of works in which similar results were obtained. A majority of these works also are not characterized by mathematical rigor, and they are based on certain physical assumptions. Therefore, one may relate them to the physical theory of diffraction. Only in a few works (related to the simpler diffraction problems) did they succeed in obtaining specific results at a higher level of mathematical rigor---more precisely, while developing asymptotic methods of mathematical diffraction theory.

We will briefly list the most important results obtained in a number of papers and books, grouping the material in the following sequence:

1. Diffraction by plane, infinitely thin plates (an infinite strip, a circular disk) and diffraction by auxiliary apertures in a flat screen (an infinite slit, a circular hole).

2. Diffraction by three-dimensional bodies with edges (a finite cylinder, a finite cone, etc.).

3. Other diffraction problems.

When investigating the first group of diffraction problems, it is necessary to keep in mind the principle of duality [4] which enables one to easily change from a strip to a slit, from a disk to a circular hole, etc. In the literature as a rule, they preferred to investigate apertures in an infinite flat screen, whereas in our book, diffraction by a strip and a disk was studied. This approach facilitates the transition to three-dimensional bodies (see the remarks at the beginning of this chapter).

Based on the time of appearance (if we do not consider the works of Schwarzschild [15] which we talked about in the Introduction), one should first of all mention the works of Braunbek [28--30] which were devoted to the diffraction of a scalar wave by a circular hole in a flat screen. Assuming that the plane wave is incident normal to
\end{scanpage}
